\documentclass[10pt,a4paper]{article}
\usepackage[utf8]{inputenc}
\usepackage[a4paper,left=3cm,right=2cm]{geometry}
\usepackage{amsmath}
\usepackage[thmmarks, amsmath, thref]{ntheorem}
\usepackage{amsfonts}
\usepackage{amssymb}
\usepackage{fancyhdr}
\pagestyle{fancy}
\usepackage{anysize}
\usepackage{graphicx}
\usepackage{float}
\usepackage{hyperref}
\usepackage{enumitem}
\usepackage{verbatim}
\usepackage{siunitx}
\newtheorem{theorem}{Theorem}
\theoremsymbol{\ensuremath{\blacksquare}}
\newtheorem*{solution}{Solución:}

\usepackage{xcolor}

% Margenes y formalidades

\textwidth 6.25in % Margen Derecho y ancho del texto
\textheight 8.5in
\oddsidemargin 0in % Margen Izquierdo
\headheight 0in % Largo del texto
\leftmargin 1in
\parindent 0em % Salto entre la indentación (Sangría)https://www.overleaf.com/4519127669sdpptkhvsbfn
\topmargin -0.5in % Margen Superior
\parskip 0.5ex % Salto entre parrafos
\baselineskip 1pt



\lhead{
\includegraphics[scale=0.08]{figuras/logoUSM.png}} % texto izquierda de la cabecera
%\chead{}                                      % texto centro de la cabecera
\rhead{Universidad Técnica Federico Santa María\\Ingeniería Comercial } % número de página a la derecha
\renewcommand{\headrulewidth}{0.4pt} % grosor de la línea de la cabecera
\renewcommand{\footrulewidth}{0.4pt} % grosor de la línea del pie

\begin{document}
\vspace{5cm}



\begin{center}
\begin{large}
\textbf{\\CERTAMEN 3}
\end{large}\\
\end{center}

\begin{center}
\begin{large}
\textbf{Econometría ICS-294}
\end{large}\\
\vspace{0.2cm}
I Semestre 2025\\
Profesores: Jocelyn Tapia Stefanoni, Franco Antonucci Fernández y Ricardo Pruneda Urz\'ua\\
%Ayudante: 
\end{center}

\textbf{Instrucciones:} 

\begin{itemize}

    \item Usted dispone de 90 minutos  para intentar los 100 puntos del certamen.
    \item Cuide la presentación y redacción de sus preguntas.
    \item No puede utilizar celular, computador o tablet, pero sí una calculadora TI.
     
\end{itemize}

\section{RdA3.2  Preguntas cortas (40 puntos)}
{\bf Elija solo 4 preguntas y responda. }
\begin{enumerate}

\item Suponga el siguiente modelo poblacional:
$$PAT = \gamma_0 + \gamma_1 \cdot \text{hombre} + \gamma_2 \cdot econ + \gamma_3 \cdot \text{hombre} \times econ +u$$

Donde: $PAT$ es el patrimonio que acumula una persona; hombre 1 si individuo es hombre, 0 si mujer; $econ = 1$ si estudia economía, 0 si no.
\begin{enumerate}
    \item[(a)] (5 puntos) ¿Cuál es el efecto esperado de haber estudiado economía sobre el patrimonio acumulado?
    {\color{blue}
    \textbf{Respuesta:} Si calculamos el efecto marginal de estudiar economía:
    $$\frac{\Delta PAT}{\Delta econ} = \gamma_2 + \gamma_3 \cdot \text{hombre}$$


En ese sentido:

\begin{itemize}
    \item $\gamma_2:$ es el efecto de estudiar economía para las mujeres.
    \item $\gamma_2 + \gamma_3:$ es el efecto de estudiar economía para los hombres.
    \item $\gamma_3$ es la diferencia del premio por estudiar economía para hombres y mujeres.
    \end{itemize}}
    \item[(b)] (5 puntos) ¿Difiere este efecto para
hombres y mujeres? Explique.

  {\color{blue} Si testeamos $\gamma_3$ y es significativo, entonces el efecto de estudiar economía es distinto para hombres y mujeres.}
\end{enumerate}








\item (10 puntos) 
Se presenta el siguiente modelo econométrico para explicar las importaciones (IMP) en función del Producto Interno Bruto (PIB), la Inversión (INV) y el Consumo (CONS):

\[
\text{IMP} = -19.7251 + 0.0322 \cdot \text{PIB} + 0.4142 \cdot \text{INV} + 0.2427 \cdot \text{CONS}
\]

Considerando la información disponible, responda lo siguiente:

\begin{enumerate}
    \item ¿Existe evidencia de multicolinealidad entre las variables explicativas? Fundamente su respuesta.

Para su respuesta considere la siguiente información.

    Matriz de correlaciones:

    \[
    \begin{array}{c|ccc}
    & \text{PIB} & \text{INV} & \text{CONS} \\
    \hline
    \text{PIB} & 1 & 0.2154 & 0.9989 \\
    \text{INV} & 0.2154 & 1 & 0.2137 \\
    \text{CONS} & 0.9989 & 0.2137 & 1 \\
    \end{array}
    \]

    Coeficientes de determinación de regresiones auxiliares para cada regresor:

    \[
    R_{PIB}^2 = 0.9979, \quad R_{INV}^2 = 0.0475, \quad R_{CONS}^2 = 0.9979
    \]

{\color{blue}
\textbf{Sí, existe evidencia clara de multicolinealidad, particularmente entre las variables \textit{PIB} y \textit{CONS}.} Esta conclusión se fundamenta en (BASTA QUE EXPLIQUEN 1):

\begin{itemize}
    \item \textbf{Matriz de correlaciones}

    La correlación entre \textit{PIB} y \textit{CONS} es $0{,}9989$, lo que indica una relación casi perfecta entre ambas variables.

    Las correlaciones entre \textit{INV} y \textit{PIB} ($0{,}2154$) e \textit{INV} y \textit{CONS} ($0{,}2137$) son bajas, lo que sugiere que \textit{INV} no presenta multicolinealidad severa.

    \item \textbf{Coeficientes de determinación de las regresiones auxiliares}

    El coeficiente de determinación $R^2$ de las regresiones auxiliares mide cuán bien una variable explicativa puede ser explicada por las otras:

    \[
    R^2_{\text{PIB}} = 0{,}9979, \quad R^2_{\text{CONS}} = 0{,}9979, \quad R^2_{\text{INV}} = 0{,}0475
    \]

    Un valor de $R^2$ cercano a 1 en las regresiones auxiliares de \textit{PIB} y \textit{CONS} indica que ambas variables están casi completamente explicadas por las otras variables del modelo.

    En particular, el alto $R^2$ sugiere que \textit{PIB} y \textit{CONS} están prácticamente colineales, lo que puede causar inestabilidad en las estimaciones del modelo.

    \item \textbf{Criterio del Factor de Inflación de la Varianza (VIF)} 

    Recordando que:

    \[
    \text{VIF} = \frac{1}{1 - R_j^2}
    \]

    Para \textit{PIB} y \textit{CONS}:

    \[
    \text{VIF} \approx \frac{1}{1 - 0{,}9979} \approx \frac{1}{0{,}0021} \approx 476{,}2
    \]

    Un VIF mayor a 10 ya es señal de multicolinealidad severa. En este caso, el VIF está sobre 400, lo cual confirma una multicolinealidad extremadamente alta.
\end{itemize}

{\color{black}
    \item En caso de confirmarse la existencia de multicolinealidad, ¿qué métodos o estrategias podrían emplearse para mitigar este problema en el modelo?
\end{enumerate}

(BASTA QUE EXPLIQUEN 1)
{\color{blue}
\begin{itemize}
    \item \textbf{Eliminar una de las variables altamente correlacionadas}
    \begin{itemize}
        \item Dado que \textit{PIB} y \textit{CONS} están altamente correlacionadas, se puede considerar eliminar una de las dos (por ejemplo, \textit{CONS} si está incluido implícitamente en el \textit{PIB}).
        \item Esto puede mejorar la precisión de las estimaciones sin perder poder explicativo.
    \end{itemize}

    \item \textbf{Combinar variables en un índice}
    \begin{itemize}
        \item Se puede construir un índice compuesto que resuma la información común entre \textit{PIB} y \textit{CONS}.
        \item Esto reduce la colinealidad al incorporar solo una nueva variable no correlacionada con las demás.
    \end{itemize}

    \item \textbf{Reespecificación del modelo}
    \begin{itemize}
        \item Puede implicar usar tasas de crecimiento o transformaciones (como logaritmos) que reduzcan la correlación entre variables.
    \end{itemize}

    \item \textbf{Recolección de más datos}
    \begin{itemize}
        \item Si es posible, aumentar el tamaño de la muestra puede ayudar a mejorar la precisión de las estimaciones, aunque no elimina la multicolinealidad.
    \end{itemize}
\end{itemize}
{\color{black}
\item (10 puntos)
¿Cómo podría utilizar el test de White (simplificado) para detectar heterocedasticidad en el modelo \( y_i = \beta_0 + \beta_1 x_{1i} + \beta_2 x_{2i} + u_i \)? Describa paso a paso el procedimiento e indique cómo se interpreta el resultado

{\color{blue}
\begin{enumerate}
    \item \textbf{Estimar el modelo original por MCO}

    \[
    y_i = \hat{\beta}_0 + \hat{\beta}_1 x_{1i} + \hat{\beta}_2 x_{2i} + \hat{u}_i
    \]

    Obtienes los residuos estimados \( \hat{u}_i \).

    \item \textbf{Elevar al cuadrado los residuos}

    Calcula:

    \[
    \hat{u}_i^2
    \]

    Esta será la variable dependiente en la regresión auxiliar.

    \item \textbf{Regresión auxiliar simplificada}

    Estima la siguiente regresión:

    \[
    \hat{u}_i^2 = \alpha_0 + \alpha_1 \hat{y}_{i}^2 + \alpha_2 \hat{y}^2_{i} + e_i
    \]

    \textit{Nota:} No se incluyen los términos originales ni los productos cruzados, sólo los cuadrados de los valores predichos de y.

    \item \textbf{Obtener el estadístico LM}

    Calcula el estadístico:

    \[
    LM = n \cdot R^2_{\text{aux}}\sim X^2_{df=2}
    \]

    Donde:

    \begin{itemize}
        \item \( n \): número de observaciones,
        \item \( R^2_{\text{aux}} \): coeficiente de determinación de la regresión auxiliar.

         El estadístico \( LM \) sigue una distribución \( \chi^2 \) con 2 grados de libertad.
    \end{itemize}

    \item \textbf{Hipótesis y decisión}

\[
H_0: \alpha_1=\alpha_2 = 0
\]
      
\[
 \alpha_1\neq 0 \quad  o \quad\alpha_2 \neq 0
\]


   

    \textbf{Decisión:}
    
    \begin{itemize}
        \item Si \( LM > \chi^2_{\alpha, 2} \), se \textbf{rechaza} \( H_0 \) → hay heterocedasticidad.
        \item Si \( LM \leq \chi^2_{\alpha, 2} \), \textbf{no se rechaza} \( H_0 \) → no hay evidencia de heterocedasticidad.
    \end{itemize}
\end{enumerate}

{\color{black}
\item  Sea el proceso estocástico \( x_t \) definido por:
\[
y_t = y_{t-1} + \varepsilon_t,
\]
donde \( \varepsilon_t \) es un proceso de ruido blanco. 

\begin{enumerate}
    \item[(a)] (3 puntos) ¿Puede afirmarse que \( y_t \) es un \textit{random walk}? Justifique conceptualmente.

    \textcolor{blue}{
        Sí, \( y_t \) es un \textit{random walk} (caminata aleatoria), porque su valor en cada instante de tiempo se construye sumando un término aleatorio (\( \varepsilon_t \)) a su valor inmediatamente anterior. Este tipo de procesos corresponde a un modelo ARIMA(0,1,0). Es una  un \textit{random walk} sin deriva.
        }
    \item[(b)]  (3 puntos) ¿Es el proceso \( y_t \) estacionario? Justifique su respuesta, considerando tanto la media como la varianza del proceso.

    \textcolor{blue}{
        La varianza del proceso crece indefinidamente, por tanto, no cumple la condición de varianza constante requerida para la estacionariedad.
        }
    \item[(c)]  (4 puntos) ¿Influye el hecho de que \( \varepsilon_t \) sea gaussiano en la clasificación de \( y_t \) como \textit{random walk}? Explique brevemente.

     \textcolor{blue}{
        No, el hecho de que \( \varepsilon_t \) sea gaussiano no influye en que \( y_t \) sea clasificado como un \textit{random walk}. La clasificación como caminata aleatoria depende de la estructura del modelo (es decir, que cada valor dependa del anterior más un término aleatorio), no de la distribución específica del ruido blanco. Sin embargo, si \( \varepsilon_t \) es gaussiano, entonces \( y_t \) será también un proceso gaussiano, lo cual puede ser relevante para ciertos análisis estadísticos.
        }
\end{enumerate}

\item (10 puntos)
En el estudio de Angrist y Krueger (1991), se utiliza el trimestre de nacimiento como variable instrumental para los años de educación en una regresión sobre ingresos. 
 Explique por qué el trimestre de nacimiento podría ser considerado un instrumento válido. 
    Explique qué es el efecto badana.

{\color{blue}  Un instrumento es v\'alido si cumple con \textbf{dos} supuestos clave:
    \begin{enumerate}
        \item {\color{blue}\textbf{Relevancia:} }$\operatorname{Cov}(X, Z) \neq 0$. La variable $Z$ afecta a $X$. Como observamos $Z$ y $X$, este supuesto es testeable.
        \item  {\color{blue}\textbf{Exogeneidad o restricci\'on de exclusi\'on:}} $\operatorname{Cov}(Z, U)=0$. 
        \begin{itemize}
            \item El instrumento no se correlaciona con el error.
            \item $E[Y \mid X]=E[Y \mid X, Z]$: el instrumento $Z$ no afecta directamente a $Y$, s\'olo lo hace a trav\'es de su efecto en $X$.
        \end{itemize}
    \end{enumerate}
 \textbf{ Validez del instrumento: trimestre de nacimiento}
    \begin{itemize}
        \item \textbf{(3 pts) Relevancia:}
        \begin{itemize}
            \item El trimestre de nacimiento está correlacionado con los años de escolaridad, ya que determina la edad de inicio escolar y, por tanto, la edad mínima legal para abandonar la escuela. \hfill (1 pt)
            \item Esta relación se confirma empíricamente al mostrar que el instrumento tiene un efecto significativo sobre la variable endógena (años de educación). \hfill (1 pt)
        \end{itemize}
        
        \item \textbf{(3 pts) Exogeneidad:}
        \begin{itemize}
            \item El trimestre de nacimiento no afecta directamente los ingresos, solo a través de su efecto sobre la educación. \hfill (1 pt)
            \item La distribución del trimestre de nacimiento es prácticamente aleatoria, por lo que no está correlacionada con habilidades no observadas u otros factores que afecten los ingresos. \hfill (1 pt)
            \item Puede haber una breve discusión sobre posibles amenazas a la exogeneidad, como efectos estacionales vinculados a decisiones familiares o salud. \hfill (1 pt)
        \end{itemize}

     Aunque la variable “años de educación” aumenta de forma continua, los títulos funcionan como puntos discretos que generan un salto en los ingresos.
El solo hecho de obtener el diploma tiene un valor adicional, incluso si los años de estudio no cambian mucho.   Se llama efecto badana por el material con el que historicamente se fabricaban los certificados de estudios.
    \end{itemize}
   }

 \item Un investigador ha simulado tres series temporales de 200 observaciones. Dos de ellas corresponden a procesos \textbf{AR(1)} con distintos valores del parámetro \( \phi \), y la tercera corresponde a un proceso de \textbf{ruido blanco} (white noise). Para cada caso se presentan dos gráficos:

\begin{itemize}
    \item A la izquierda: la serie temporal simulada \( \{X_t\} \)
    \item A la derecha: su función de autocorrelación (ACF) hasta el rezago 30
\end{itemize}

Los valores de \( \phi \) utilizados para los procesos AR(1) fueron:

\[
\phi = -0.95 \quad \text{y} \quad \phi = 0.95
\]

\[
X_t = \phi X_{t-1} + \varepsilon_t
\]

Los gráficos han sido rotulados como \textbf{Caso A}, \textbf{Caso B} y \textbf{Caso C} (sin indicar el valor de \( \phi \) ni el tipo de proceso que representan).

\begin{enumerate}
    \item (4 puntos)  Identifique cuál de los tres casos corresponde al proceso de \textbf{ruido blanco}. Justifique su respuesta con base en el comportamiento de la serie y su ACF.


\textcolor{blue}{
        El \textbf{Caso B} corresponde al proceso de ruido blanco. La serie temporal no muestra una estructura clara ni dependencia entre observaciones sucesivas. Además, (ANALISIS MÁS IMPORTANTE PUEDE TENER TODO EL PUNTAJE CON ESTA PARTE) la función de autocorrelación (ACF) muestra valores cercanos a cero para todos los rezagos, lo cual es característico de un proceso de ruido blanco.
        }

    \item (6 puntos) Para los otros dos casos, asigne a cada uno el valor de \( \phi \) correspondiente. Justifique su elección explicando cómo el signo y la magnitud de \( \phi \) se reflejan en la dinámica de la serie y en la forma de su ACF.

{\color{blue}
        El \textbf{Caso A} corresponde al proceso AR(1) con \( \phi = -0.95 \). Se observa un patrón oscilante en la serie, donde los valores tienden a alternar de signo. Esta dinámica es típica de un proceso AR con un coeficiente negativo cercano a -1. La ACF también muestra un patrón alternado (signo negativo en rezagos impares) (ANALISIS MÁS IMPORTANTE PUEDE TENER TODO EL PUNTAJE CON ESTA PARTE) , decreciendo gradualmente, lo que confirma el signo negativo de \( \phi \).

        El \textbf{Caso C} corresponde al proceso AR(1) con \( \phi = 0.95 \). La serie muestra una fuerte persistencia: los valores permanecen en la misma dirección por varios periodos antes de cambiar, lo cual indica una alta autocorrelación positiva. La ACF (ANALISIS MÁS IMPORTANTE PUEDE TENER TODO EL PUNTAJE CON ESTA PARTE)  decrece lentamente y de forma positiva, lo que es consistente con un valor de \( \phi \) positivo cercano a 1.
        }   
\end{enumerate}

\begin{figure}[H]
    \centering
    \includegraphics[width=0.7\linewidth]{certamenes/autocorr.png}
\end{figure}

\end{enumerate}


\section{RdA3.2  Ejercicios (60 puntos)}


{\bf A continuación elija solo 3 preguntas y responda. }

\begin{enumerate}
 \item  Con ayuda de los datos del archivo \texttt{GPA2.RAW} se obtuvo la siguiente ecuación estimada:

\[
\widehat{sat} =
\underset{(6.29)}{1028.10} +
\underset{(3.83)}{19.30} \cdot hsize -
\underset{(0.53)}{2.19} \cdot hsize^2 +
\underset{(4.29)}{45.09} \cdot female -
\underset{(12.71)}{169.81} \cdot lowincome +
\underset{(18.15)}{62.31} \cdot (female \times lowincome)
\]


\vspace{1em}
\noindent
Otros datos:
\[
n = 4{,}137, \quad R^2 = 0.0858
\]

\vspace{1em}
\noindent
Las variables utilizadas son:

\begin{itemize}
    \item \textbf{sat:} Puntuación combinada en el examen de admisión (SAT).
    \item \textbf{hsize:} Tamaño de la cohorte de estudiantes (en centenas) que terminaron el bachillerato en la escuela del estudiante.
    \item \textbf{female:} Variable binaria que vale 1 si el estudiante es mujer y 0 si es hombre.
    \item \textbf{lowincome:} Variable binaria que vale 1 si el estudiante es de bajos ingresos y 0 en caso contrario.
    \item \textbf{female $\times$ lowincome:} Interacción entre género e ingreso bajo.
\end{itemize}

\vspace{1em}

\begin{enumerate}
    \item[i)] (5 puntos) ¿Existe alguna evidencia fuerte de que $hsize^2$ deba ser incluida en el modelo? Con base en esta ecuación, ¿cuál es el tamaño óptimo de una escuela?

{\color{blue}
\textbf{Hipótesis para \( hsize^2 \):}

\[
H_0: \beta_{hsize^2} = 0 \quad vs \quad H_1: \beta_{hsize^2} \neq 0
\]

Coeficiente estimado:  
\[
\hat{\beta}_{hsize^2} = -2.19
\]

Error estándar:  
\[
SE = 0.53
\]

Estadístico \( t \):

\[
t = \frac{-2.19}{0.53} \approx -4.13
\]

El valor absoluto es mayor que 1.96 (nivel 5\%), por lo tanto:

\[
\boxed{
\text{Rechazamos } H_0 \Rightarrow \text{Existe evidencia significativa para incluir } hsize^2
}
\]

\vspace{1em}

\textbf{Tamaño óptimo de la escuela:} es donde la puntuación estimada del SAT se maximiza respecto a \( hsize \):

\[
\frac{d \widehat{sat}}{d hsize} = 19.30 - 2 \times 2.19 \times hsize = 0
\]

\[
19.30 - 4.38 \cdot hsize = 0 \quad \Rightarrow \quad hsize^* = \frac{19.30}{4.38} \approx 4.41
\]

\textbf{Interpretación:}  
El tamaño óptimo de la escuela es aproximadamente \(4.41\) centenas de estudiantes, es decir, \(\boxed{441}\) estudiantes.

{\color{black}    
    \item[ii)] (5 puntos) Manteniendo $hsize$ constante, ¿cuál es la diferencia estimada en la puntuación del SAT entre mujeres de ingresos altos y hombres de altos ingresos? ¿Es estadísticamente significativa  esta diferencia?

{\color{blue}  
Para ingresos altos:  
\[
lowincome = 0
\]

Entonces, la diferencia es sólo por el coeficiente de \( female \):

\[
\Delta = \hat{\beta}_{female} = 45.09
\]

Error estándar:  
\[
SE = 4.29
\]

Estadístico \( t \):

\[
t = \frac{45.09}{4.29} \approx 10.51
\]

\textbf{Conclusión:}

Las mujeres de ingresos altos tienen en promedio \(45.09\) puntos más que los hombres de ingresos altos.

La diferencia es estadísticamente significativa al nivel \(5\%\).

{\color{black}      
    \item[iii)] (5 puntos) ¿Cuál es la diferencia estimada en la puntuación del SAT entre hombres de ingresos altos y hombres de ingreso bajo? Pruebe la hipótesis nula de que no hay diferencia entre estas puntuaciones contra la alternativa de que sí hay diferencia.
    
{\color{blue} 
Para hombres:  
\[
female = 0
\]

La diferencia viene del coeficiente \( lowincome \):

\[
\Delta = \hat{\beta}_{lowincome} = -169.81
\]

Error estándar:  
\[
SE = 12.71
\]

Estadístico \( t \):

\[
t = \frac{-169.81}{12.71} \approx -13.36
\]

\textbf{Prueba:}

\[
H_0: \beta_{lowincome} = 0 \quad vs \quad H_1: \beta_{lowincome} \neq 0
\]


Existe suficiente evidencia estadística para Rechazar Ho con un nivel de significancia del 5\%. Considerando que la muestra tiene más de 4000 datos usamos los valores críticos de la normal.  }



  {\color{black}   
    \item[iv)] (5 puntos) ¿Cuál es la diferencia estimada en la puntuación del SAT entre mujeres de bajos ingresos y mujeres de ingresos altos? ¿Qué se necesita hacer para probar si esta diferencia es estadísticamente significativa?

{\color{blue} 
Mujeres:  
\[
female = 1
\]

Diferencia:

\[
\Delta = \hat{\beta}_{lowincome} + \hat{\beta}_{female \times lowincome} = -169.81 + 62.31 = -107.5
\]

Para evaluar si esta diferencia es estadísticamente significativa. Basta con chequear que al menos uno de los coeficientes es significativos. 

\[t_{\beta_{lowincome}} ={-169,81\over 12,71}=-13,36\]

\[t_{\beta_{female \times lowincome}} ={62,31\over 18,15}=3,43\]

Considerando que la muestra tiene más de 4000 observaciones utilizamos los valores críticos de la normal y por tanto ambos son significativos al 1\%.
}
\end{enumerate}
  \item Un profesor está analizando qué factores explican el rendimiento académico de los estudiantes en su curso de Econometría. Estima el siguiente modelo de regresión:

\[
\text{nota}_i = \beta_0 + \beta_1 \cdot \text{estudio}_i + \beta_2 \cdot \text{trabaja}_i + u_i
\]

Donde:
\begin{itemize}
    \item $\text{nota}_i$: nota final del estudiante \(i\) en el curso,
    \item $\text{estudio}_i$: horas promedio de estudio a la semana,
    \item $\text{trabaja}_i$: variable dicotómica que vale 1 si el estudiante trabaja durante el semestre, 0 en caso contrario,
    \item $u_i$: término de error, que incluye factores no observados como concentración, salud, o estrés.
\end{itemize}

El profesor estima su modelo y obtiene lo siguiente: %Después de estimar el modelo, el profesor nota que los \textbf{residuos tienden a ser mayores (en magnitud) para estudiantes que trabajan}. Sospecha entonces que hay \textbf{heterocedasticidad} en el modelo.

\section*{Resultados de R\footnote{  ** indica 5\% de significancia, *** indica 1\% de significancia}}

\noindent\fbox{%
\begin{minipage}{0.95\textwidth}
\ttfamily
\small
Call:\newline
lm(formula = nota ~ estudio + trabaja, data = datos)\newline

Residuals:\newline
\hspace*{1em}Min \hspace{3em} 1Q \hspace{3em} Median \hspace{2em} 3Q \hspace{3em} Max\newline
\hspace*{1em}-6.72 \hspace{1.5em} -1.63 \hspace{2em} -0.04 \hspace{3em} 1.46 \hspace{3em} 5.75\newline

Coefficients:\newline
\hspace*{1em}Estimate \hspace{1em} Std. Error \hspace{1em} t value \hspace{1em} Pr(>|t|)\newline
(Intercept) \hspace{0.6em} 4.3262 \hspace{2em} 1.6280 \hspace{2em} 2.657 \hspace{2em} 0.00901 **\newline
estudio \hspace{2.3em} 0.4589 \hspace{2em} 0.1630 \hspace{2em} 2.818 \hspace{2em} 0.00587 **\newline
trabaja \hspace{2.3em} -2.8665 \hspace{1.9em} 0.5962 \hspace{2em} -4.807 \hspace{1.7em} 0.00001 ***\newline

Residual standard error: 2.91 on 97 degrees of freedom\newline
Multiple R-squared: 0.235,\quad Adjusted R-squared: 0.219\newline
F-statistic: 14.87 on 2 and 97 DF, \quad p-value: 2.32e-06
\end{minipage}
}


\begin{enumerate}
    \item (5 puntos) ¿Por qué podría haber heterocedasticidad en este contexto? ¿Por qué representa un problema para la estimación por MCO?\\
{\color{blue}
    \textbf{Respuesta:} La heterocedasticidad implica que la varianza del término de error no es constante entre observaciones. En este caso, los estudiantes que trabajan presentan mayor variabilidad en su rendimiento. Aunque los estimadores MCO siguen siendo insesgados, sus errores estándar serán incorrectos, lo que invalida los tests de significancia e intervalos de confianza.}\\

    \item (5 puntos) Plantee la regresión auxiliar del test de White ampliado para este problema.\\

%\begin{comment}
{\color{blue}
    \textbf{Respuesta:} La regresión quedaría de la forma:
$$\hat{u}_i^2 = \gamma_0 
+ \gamma_1 \cdot \text{estudio}_i 
+ \gamma_2 \cdot \text{trabaja}_i 
+ \gamma_3 \cdot \text{estudio}_i^2 
+ \gamma_4 \cdot \text{trabaja}_i^2 
+ \gamma_5 \cdot (\text{estudio}_i \cdot \text{trabaja}_i) 
+ v_i$$}
%\end{comment}

    \item (10 puntos) Plantee el test de hipótesis de heterocedasticidad. Si el $R^2$ de la regresión auxiliar asciende al valor de 25\%, ¿existe o no heterocedasticidad al 5\% de significancia? Justifique su respuesta.\\


%\begin{comment}
{\color{blue}
\textbf{Respuesta:}

El test de White evalúa la presencia de heterocedasticidad a través de la regresión auxiliar sobre los residuos al cuadrado.

\begin{itemize}
    \item (\(H_0\)): Homocedasticidad, es decir, la varianza de los errores es constante. Matemáticamente: \(\operatorname{Var}(u_i) = \sigma^2\).
    \item (\(H_1\)): Heterocedasticidad, es decir, la varianza de los errores depende de una o más variables explicativas.
\end{itemize}

El estadístico de prueba es:

\[
n \cdot R^2 \sim \chi^2_{k}
\]

donde:
\begin{itemize}
    \item \(n\) es el número de observaciones (en este caso, 100),
    \item \(R^2 = 0{,}25\) es el coeficiente de determinación de la regresión auxiliar,
    \item \(k\) es el número de regresores (sin incluir la constante). Para el test de White ampliado con dos regresores originales, se tienen: \texttt{estudio}, \texttt{trabaja}, \texttt{estudio\textsuperscript{2}}, \texttt{estudio × trabaja}, \texttt{trabaja\textsuperscript{2}}, es decir, \(k = 5\).
\end{itemize}

Entonces:

\[
n \cdot R^2 = 100 \cdot 0.25 = 25
\]

Se compara con el valor crítico de la distribución \(\chi^2_5\) al 5\% de significancia:

\[
\chi^2_{5,\,0.95} \approx 11.07
\]

Como:

\[
25 > 11.07
\]

Se \textbf{rechaza la hipótesis nula} de homocedasticidad. Por lo tanto, existe evidencia estadísticamente significativa de \textbf{heterocedasticidad} en el modelo.
}
%\end{comment}

\end{enumerate}

\item A continuación se presenta el gráfico de la serie temporal del tipo de cambio nominal (pesos por dólar) en Chile, correspondiente al período comprendido entre enero de 1983 y diciembre de 2024, con datos de frecuencia mensual. 


\begin{center}
\textbf{Figura 1: Tipos de cambio (pesos por dólar)}
\begin{figure}[H]
    \centering
    \includegraphics[width=0.85\linewidth]{certamenes/BCChChart.png}
\end{figure}
\end{center}

Con el objetivo de evaluar si la serie temporal sigue un proceso de caminata aleatoria (random walk), un grupo de investigadores aplica el test de Dickey-Fuller mediante la estimación de dos modelos auxiliares:

   \textbf{Modelo sin tendencia (without trend):}
    \begin{table}[H]
    \centering
    \begin{tabular}{lccc}
    \toprule
    Variable & Coef. & Error estándar & Estadístico t \\
    \midrule
    Intercepto & 0,0710 & 0,0354 & 2,006 \\
    $y_{t-1}$  & -0,0182 & 0,0097 & -1,876 \\
    \bottomrule
    \end{tabular}
    \end{table}

    \textbf{Modelo auxiliar con tendencia (\textit{trend}):}
    \begin{table}[H]
    \centering
    \begin{tabular}{lccc}
    \toprule
    Variable & Coef. & Error estándar & Estadístico t \\
    \midrule
    Intercepto & 0,0452 & 0,0289 & 1,563 \\
    Tendencia ($t$) & 0,0021 & 0,0003 & 7,000 \\
    $y_{t-1}$  & -0,0304 & 0,018 & -1,6889 \\
    \bottomrule
    \end{tabular}
    \end{table}

  Testee con un 1\% de significancia si el tipo cambio nominal tiene raíz unitaria. Explique su resultado.  


{\color{blue}
  
La inspección visual de la serie permite identificar a un crecimiento sistemático a lo largo del tiempo, compatible con una tendencia temporal. Por tanto, se justifica el uso del test Dickey-Fuller con tendencia.


Modelo auxiliar: \\
\[ 
\Delta y_t = \alpha + \beta t + \phi y_{t-1} + \varepsilon_t 
\]
Hipótesis:
\begin{itemize}
    \item \( H_0: \phi = 0 \) → la serie tiene raíz unitaria (no estacionaria).
    \item \( H_1: \phi < 0 \) → la serie es estacionaria alrededor de una tendencia determinista.
\end{itemize}


Si se rechaza \( H_0 \), se concluye que el proceso es estacionario alrededor de una tendencia determinista, es decir, fluctúa de forma estable en torno a una línea recta en el tiempo.


\begin{itemize}
    \item Sin tendencia: Estadístico t para \( y_{t-1} \) = $-1{,}876$
    \item Con tendencia: Estadístico t para \( y_{t-1} \) = $-1{,}6889$
\end{itemize}


Se utiliza el valor del estadístico t del modelo con tendencia: $-1{,}6889$.\\

La serie tiene una frecuencia mensual desde enero de 1983 a diciembre de 2024  con un total de 504 observaciones.

Se deben utilizar los valores críticos de la tabla para T=500:

\begin{itemize}
    \item 1\%: $-3{,}98$
    \item 5\%: $-3{,}42$
\end{itemize}

Comparación: \(-1{,}6889 > -3{,}98\) → No se rechaza \( H_0 \)  al 1\%.\\

La evidencia no permite rechazar la presencia de raíz unitaria. La serie es no estacionaria y sigue una caminata aleatoria con tendencia estocástica.
}

 \item Un investigador desea estudiar el efecto de las \textbf{horas de entrenamiento semanal} en el \textbf{rendimiento deportivo} de jugadores de básquetbol universitario. Plantea el siguiente modelo:

 
\[
y_i = \beta_0 + \beta_1 \cdot x_i + \beta_2 \cdot h_i + u_i
\]

Donde:
\begin{itemize}
    \item $y_i$: promedio de puntos anotados por partido durante la temporada.
    \item $x_i$: número de horas de entrenamiento individual por semana. 
    \item $h_i$: altura del jugador en centímetros (variable de control). 
    \item $u_i$: factores no observados que afectan el rendimiento (motivación, talento, disciplina, etc.).
\end{itemize}



Sin embargo, se sospecha que la cantidad de horas de entrenamiento está correlacionada con $u_i$. Por ejemplo, los jugadores más motivados podrían entrenar más y también rendir mejor, incluso si el entrenamiento no fuera la causa del mayor rendimiento.


\begin{enumerate}
    \item (5 puntos) ¿Por qué podría ser inválido estimar el modelo por Mínimos Cuadrados Ordinarios (MCO)? Explique con base en los supuestos del modelo clásico.\\
{\color{blue}
    \textbf{Respuesta:} MCO será inválido si $\text{entreno}_i$ está correlacionado con el término de error $u_i$, violando el supuesto clásico de exogeneidad (\textit{E}[$u_i | X_i$] = 0). Esto ocurre, por ejemplo, si jugadores más motivados entrenan más y además tienen un rendimiento mayor por razones distintas al entrenamiento.}

    \item (5 puntos) Explique por qué $\text{distancia\_cancha}_i$ medida como ``la distancia (en minutos caminando) entre el dormitorio del jugador y la cancha de entrenamiento'' (le llamaremos z) podría ser un \textbf{instrumento válido} para $\text{entreno}_i$. ¿Qué condiciones debe cumplir z para ser un buen instrumento? Exprese matemáticamente.\\
%\begin{comment}
{\color{blue}
    \textbf{Respuesta:} $\text{distancia\_cancha}_i$ podría ser un instrumento válido si cumple:
    \begin{itemize}
        \item \textbf{Relevancia:} Está correlacionada con $\text{entreno}_i$. En este caso, mayor distancia reduce la probabilidad de entrenar más $\implies Cov(x_i, z_i) \neq 0$.
        \item \textbf{Exogeneidad:} No está correlacionada con $u_i$; es decir, no afecta directamente el rendimiento, ni está relacionada con factores como la motivación o el talento $\implies Cov(u_i, z_i) = 0$.
    \end{itemize}
    }
%\end{comment}

    \item (5 puntos) Suponga que $\operatorname{Cov}(x_i, z_i) = -2,94765$ y $\operatorname{Cov}(y_i, z_i) -4,3215$ ¿Qué valor debería tener $\beta_1$? Derive la operación matemática e interprete.\\

%\begin{comment}
{\color{blue}
   \textbf{Respuesta:} El valor de $\beta_1$ se puede obtener a partir de la covarianza entre la variable dependiente $y_i$ (puntos) y el instrumento $z_i$ (distancia):

\[
\begin{aligned}
\operatorname{Cov}(y_i, z_i) &= \operatorname{Cov}(\beta_0 + \beta_1 x_i + \beta_2 h_i + u_i, z_i) \\
&= \underbrace{\operatorname{Cov}(\beta_0, z_i)}_{= 0} + \beta_1 \operatorname{Cov}(x_i, z_i) + \underbrace{\operatorname{Cov}(h_i, z_i)}_{= 0} + \underbrace{\operatorname{Cov}(u_i, z_i)}_{= 0} \\
&= \beta_1 \operatorname{Cov}(x_i, z_i)
\end{aligned}
\]

Despejando:

\[
\beta_1 = \frac{\operatorname{Cov}(y_i, z_i)}{\operatorname{Cov}(x_i, z_i)} \implies \beta_1 = \frac{-4,3215}{-2,94765} = 1,4661
\]

Como la distancia a la cancha reduce las horas de entrenamiento, y entrenar más aumenta el rendimiento, se espera que $\operatorname{Cov}(x_i, z_i) < 0$ y $\operatorname{Cov}(y_i, z_i) < 0$, por lo que $\beta_1 > 0$: entrenar más tiene un efecto positivo sobre los puntos anotados, particularmente, una hora más de entrenamiento a la semana genera que el promedio de puntos suba 1,47 puntos.\\
    }
%\end{comment}

    \item (5 puntos) ¿El coeficiente estimado mediante variable instrumental será mayor o menor que el estimado por MCO?\\
{\color{blue}
    \textbf{Respuesta:} Si la motivación no observada está positivamente correlacionada con las horas de entrenamiento, entonces el estimador MCO de $\beta_1$ estará sesgado hacia arriba: capta tanto el efecto del entrenamiento como el de la motivación. En cambio, la estimación por variable instrumental elimina este sesgo. Por tanto, el coeficiente estimado por IV será \textbf{menor} que el estimado por MCO.}

\end{enumerate}

\item Suponga que usted está interesado en determinar los factores que explican
las diferencias salariales entre distintos individuos. Para realizar
su análisis, usted especifica la siguiente función de regresión:

\begin{equation}
wage = \beta_{1} + \beta_{2}IQ + \beta_{3}Educ + \beta_{4}Exper + \beta_{5}Tenure + \beta_{6}personacolor + u \label{eq:original}
\end{equation}

donde \emph{wage} representa el salario, \emph{IQ} es el coeficiente intelectual,
\emph{Educ} son los años de educación, \emph{Exper} los años de experiencia laboral,
\emph{Tenure} la antigüedad en el trabajo actual, y \emph{personacolor} es una variable dicotómica
que toma el valor 1 cuando la persona se identifica como {persona de color}, y 0 en caso contrario.


Los resultados de la regresión se presentan en la Tabla 1, y en la Figura 1 se muestran los residuos
agrupados según la variable \emph{personacolor}:

\begin{center}
\textbf{Tabla 1}
\end{center}

\begin{center}
Model 1: OLS, using observations 1--935\\ Dependent variable: wage\\
\vspace{1em}
\begin{tabular}{lr@{.}lr@{.}lr@{.}lr@{.}l}
 &  \multicolumn{2}{c}{Coefficient} &   \multicolumn{2}{c}{Std.\ Error} &    \multicolumn{2}{c}{$t$-ratio} &     \multicolumn{2}{c}{p-value} \\
\hline
const &   $-$405&803 &     122&400 &       $-$3&3154 &         0&0010 \\
IQ &   3&77957 &     0&996294 &       3&7936 &         0&0002 \\
educ &   57&5438 &     7&00006 &       8&2205 &         0&0000 \\
exper &   14&8737 &     3&19330 &       4&6578 &         0&0000 \\
tenure &   6&88539 &     2&46227 &       2&7964 &         0&0053 \\
personacolor & -40&15 & 38&70 & -1&037 & 0&296\\
\end{tabular}
\vspace{1ex}

\begin{tabular}{lrlr}
Mean dependent var &  957.9455 & S.D. dependent var &  404.3608 \\
Sum squared resid &  1.25\textrm{e+08} & S.E. of regression &  367.4213 \\
$R^2$ &  0.178780 & Adjusted $R^2$ &  0.174361 \\
$F(5, 929)$ &  40.44887 & P-value($F$) &  1.07\textrm{e--37} \\
\end{tabular}
\end{center}

\begin{center}
\textbf{Figura 1: Residuos agrupados por identificación étnico-racial}
\begin{figure}[H]
    \centering
    \includegraphics[width=0.65\linewidth]{certamenes/racial.png}
\end{figure}
\end{center}

\begin{enumerate}
    \item (4 puntos) De acuerdo a esta información, ¿le parece que existe heterocedasticidad en el modelo? ¿Tiene sentido este resultado? Explique.

 

    \item (7 puntos)  En la Tabla 2 se presentan los resultados de una regresión de los residuos al cuadrado de la regresión de la tabla 1 sobre las variables explicativas:

\begin{center}
\textbf{Tabla 2: OLS, usando 935 observaciones} \\
Variable dependiente: $\hat{u}^2$ \\
\vspace{1em}
\begin{tabular}{lr@{.}lr@{.}lr@{.}lr@{.}l}
 & \multicolumn{2}{c}{Coeficiente} & \multicolumn{2}{c}{Error estándar} & \multicolumn{2}{c}{$t$-ratio} & \multicolumn{2}{c}{p-valor} \\[1ex]
const          & $-$160277&000 & 104551&00 & $-$1&5330 & 0&1256 \\
IQ             & 921&598   & 851&010  & 1&0829  & 0&2791 \\
educ           & 16288&600 & 5979&280 & 2&7242  & 0&0066 \\
exper          & 2955&490  & 2727&640 & 1&0835  & 0&2789 \\
tenure         & $-$6586&500 & 2103&220 & $-$3&1316 & 0&0018 \\
personacolor   & $-$70000&000 & 30000&000 & $-$2&3333 & 0&0200 \\
\end{tabular}

\vspace{1ex}
\begin{tabular}{lrlr}
Media variable dependiente & 134132.1 & Desv. estándar dependiente & 317721.3 \\
Suma de residuos al cuadrado & $9.15 \times 10^{13}$ & Error estándar de regresión & 313842.4 \\
$R^2$ & 0.029491 & $R^2$ ajustado & 0.024268 \\
$F(5, 929)$ & 5.646001 & Valor-p ($F$) & 0.000039 \\
\end{tabular}
\end{center}


    \[
    \hat{\mu}^2 = a_1 + a_2 IQ + a_3 Educ + a_4 Exper + a_5 Tenure + a_6 personacolor + v
    \]

    ¿Puede usarse esta regresión para testear la presencia de heterocedasticidad? Explique e interprete.

      \item (5 puntos) Observe que la variable \texttt{personacolor} no es estadísticamente significativa según los resultados de la tabla 1. ¿Podemos concluir que \texttt{personacolor} no afecta el ingreso? Fundamente su respuesta considerando la presencia de heterocedasticidad. 

        \item (4 puntos) ¿Cómo se puede corregir el problema anterior para obtener inferencias válidas? ¿Qué tipo de errores estándar se deben utilizar? 

   
\end{enumerate}

\end{enumerate}

\begin{enumerate}
    \item {\color{blue}
    A partir de la Figura 1, se observa que la dispersión de los residuos parece variar dependiendo del valor de la variable \emph{personacolor}. En particular, los residuos correspondientes a personas de color presentan una mayor dispersión en comparación con el resto de la muestra, lo cual sugiere la presencia de heterocedasticidad.

    Este resultado tiene sentido desde el punto de vista económico si se considera que las personas de color pueden enfrentar condiciones laborales más heterogéneas, lo que se traduce en una mayor variabilidad no explicada del salario.}

    \item {\color{blue}
    Esta regresión puede utilizarse para testear la presencia de heterocedasticidad utilizando el test de White. La idea central es estimar una regresión auxiliar donde los residuos al cuadrado $\hat{u}^2$ se explican en función de las variables originales del modelo.

    En este caso, se aplica el test de White a través del siguiente contraste:

    \[
    H_0: a_2 = a_3 = a_4 = a_5 = a_6 = 0 \quad \text{(homocedasticidad)}
    \]
    \[
    H_1: \text{al menos un } a_j \neq 0 \quad \text{(heterocedasticidad)}
    \]

    Se utiliza como estadístico de contraste el multiplicador de Lagrange:

    \[
    LM = n \cdot R^2 = 935 \cdot 0.029491 \approx 27.57
    \]

    Este estadístico se distribuye asintóticamente como una $\chi^2$ con 5 grados de libertad (tantos como regresores excluyendo la constante). El valor crítico al 5\% es aproximadamente 11.07.

    Dado que $27.57 > 11.07$, se rechaza la hipótesis nula de homocedasticidad. Por lo tanto, se concluye que existe evidencia estadísticamente significativa de heterocedasticidad en el modelo.}


 \item {\color{blue}
    No necesariamente. Dado que hay evidencia de heterocedasticidad, el valor-p que indica que \texttt{color} no es significativa puede ser inválido. La heterocedasticidad puede haber inflado o reducido artificialmente el error estándar de ese coeficiente. Por tanto, no es posible concluir con certeza que \texttt{color} no tiene efecto.
    }   

  \item   {\color{blue}
    Para obtener inferencias válidas en presencia de heterocedasticidad, se deben utilizar errores estándar robustos. En particular, se recomienda emplear los errores estándar de tipo White (robustos a heterocedasticidad), que corrigen la matriz de varianzas-covarianzas sin asumir homocedasticidad.
    }
\end{enumerate}


\end{enumerate}








\newpage


\begin{center}
    {\Large Formulario  }

 
\end{center}

\begin{itemize}

% \item  
% {\color{red}
% $$\hat{\beta}_0 = \bar{y} - \hat{\beta}_1 \bar{x} $$
% $$\hat{\beta}_1 = \frac{\sum_{i=1}^{n}(x_i - \bar{x})(y_i-\bar{y})}{\sum_{i=1}^{n}(x_i - \bar{x})^2}$$}

\item  {\color{red} \[ \hat{\beta} = (X^{\top}X)^{-1}X^{\top}y \]}

\item 
	{\color{red}	\[ SST = \sum_{i=1}^{n} (y_i - \bar{y})^2 \quad SSE = \sum_{i=1}^{n} (\hat{y}_i - \bar{y})^2  \quad SSR = \sum_{i=1}^{n} \hat{u}_i^2\]}
		

    \item {\color{red} \[ R^2  = 1 - \frac{SSR}{SST} \quad \quad \bar{R}^2 =  1 - \frac{(n-1)SSR}{(n-k-1)SST} \]}

 \item {\color{red} \[ t_{\hat{\beta_j}}=\frac{\hat{\beta}_j - \beta_j}{\text{s.e.}(\hat{\beta}_j)} \sim t_{n-k-1} \]}

 \item {\color{red}\[
    F = \frac{SSR_{r} - SSR_{ur}}{SSR_{ur}} \cdot \frac{n - k - 1}{j} \sim F(j, n - k - 1)
    \]}

\item {\color{red}Test F para regresión auxiliar.}
\begin{center}
 {  \color{red}     \(F=\dfrac{R^{2}/j}{(1-R^{2})/(n-k-1)}
          \sim F(j,\;n-k-1)\).}
\end{center}
% \item {\color{red}Considere que el valor crítico a un nivel de significancia del 5\% es 1.645}
% \item {\color{red}Considere los siguientes valores críticos para t: $\alpha = 0.01 \implies t \approx 2.62$; $\alpha = 0.05 \implies t \approx 1,97$;$\alpha = 0.01 \implies t \approx 2.62$; $\alpha = 0.1 \implies t \approx 1,645$}
%\item {\color{red}Considere que el valor crítico de F al 1\% con grados de libertad (2,178) es aproximadamente 4.72}


\item \color{red}{ Multiplicador de Lagrange

$$LM = nR^2 \sim \chi^2_{df}$$

Tabla $\chi^2_{df}$
\begin{center}
    \begin{tabular}{|c|c|c|c|}
    \hline
        df & 0.99 & 0.95 & 0.90 \\ 
        \hline
        1 &  6.635 & 3.841 & 2.706\\
        \hline
        2 &  9.210 & 5.991 & 4.605 \\
        \hline
        3 & 11.345 & 7.815 & 6.251  \\
        \hline
        4 &  13.277 & 9.488 & 7.779 \\
        \hline
        5 &  15.086 & 11.070 & 9.236 \\
        \hline
        6 &  16.812 & 12.592 & 10645 \\
        \hline
    \end{tabular}
    \end{center}}


\item Tabla de test D-F
\begin{center}
    \begin{figure}[H]
        \centering
        \includegraphics[width=0.7\linewidth]{sem14/tabla.png}
    \end{figure}
\end{center}

\end{itemize}


% \begin{figure}[H]
%     \centering
%     \includegraphics[width=0.85\linewidth]{Ejercicios//ejercicios_inferencia/cert2_F2_1184.png}
% \end{figure}

% \begin{figure}[H]
%     \centering
%     \includegraphics[width=0.85\linewidth]{certamenes/normal.png}
% \end{figure}


\begin{figure}
    \centering
    \includegraphics[width=0.85\linewidth]{certamenes/normal_1cola.png}
\end{figure}



\end{document}